\documentclass{beamer}

\usepackage{cuhkbeamer}

\begin{document}

% ====================== Slide 1: TITLE SLIDE (REVISED) ======================
\begin{frame}
  \title{Sustainability Challenges in Urban Air Operations}
  \author{Chiu Yu KO}
  \institute{Low Altitude Economics}
  \date{Week 10}
  \titlepage
\end{frame}

% ====================== Slide 2: FORMAT & OBJECTIVES ======================
\begin{frame}[t]
\frametitle{Learning Objectives}
By the end of this class, you should be able to:
\begin{itemize}
\item \textbf{Identify} major environmental and social externalities associated with urban air operations.
\item \textbf{Compare} pricing, operating restrictions, standards, mitigation, and compensation as policy responses.
\item \textbf{Calculate} simple operational and life-cycle greenhouse-gas indicators using transparent boundaries.
\item \textbf{Evaluate} privacy and public-acceptance measures using necessity, proportionality, and stakeholder evidence.
\item \textbf{Design} measurable sustainability KPIs and assess who benefits, who pays, and who bears residual impacts.
\end{itemize}
\end{frame}

% ====================== Slide 3: ICEBREAKER ======================
\begin{frame}[t]
\frametitle{Can an LAE Service Be Sustainable and Viable?}
\textbf{Imagine a proposed electric-aircraft corridor near a dense urban community.}
\begin{itemize}
\item The service may reduce travel time or improve delivery and public-service performance.
\item Electricity use, vehicle manufacturing, noise, visual effects, data collection, and infrastructure still create impacts.
\end{itemize}
\vspace{0.3cm}
``Electric'' does not automatically mean zero-emission, quiet, private, equitable, or publicly accepted.
\begin{itemize}
\item \textbf{Quick poll (2 min):} Which issue is most likely to constrain your project: noise, privacy, greenhouse-gas emissions, energy demand, or unequal distribution of benefits?
\end{itemize}
\end{frame}

% ====================== Session 1 Title Card ======================
\begin{frame}[c]
  \begin{center}
    \huge\color{cuhkpurple}\textbf{Session 1}\\
    \vspace{0.5cm}
    \Large Externalities, Pigovian Instruments \& Privacy Governance
  \end{center}
\end{frame}

% ====================== Slide 4: CONCEPT PRIMER 1 ======================
\begin{frame}[t]
\frametitle{Concept Primer: Externalities}
An externality arises when an activity affects third parties and those effects are not fully reflected in the decision-maker's private costs or benefits.
\begin{itemize}
\item \textbf{Noise and visual effects}: annoyance, sleep disturbance, amenity, and land-use impacts.
\item \textbf{Privacy and data}: excessive collection, surveillance concerns, security risk, and loss of trust.
\item \textbf{Climate and energy}: operational electricity emissions, supply-chain emissions, and peak-load requirements.
\item \textbf{Positive effects}: time savings, safer inspection, emergency response, knowledge, and network benefits.
\end{itemize}
\textbf{Managerial lesson:} First identify and measure the impact; then compare prevention, mitigation, pricing, compensation, and operating restrictions.
\end{frame}

% ====================== Slide 5: STANDALONE FORMULA ======================
\begin{frame}[t]
\frametitle{Applied Tool: Pricing a Noise Externality}
\begin{block}{Economic Benchmark}
\[
\text{Efficient Charge per Operation}
=\text{Marginal External Noise Cost}
\]
\end{block}
\begin{itemize}
\item Marginal external cost depends on exposure, time, location, number of affected people, baseline noise, event frequency, and health or amenity effects.
\item Decibels are logarithmic; ``decibels above a limit'' should not automatically be multiplied by flights and a fixed monetary rate.
\item Practical instruments may use certified noise categories, time-of-day charges, quotas, curfews, route restrictions, or performance standards.
\item Charges should encourage quieter operations without implying that every impact can simply be paid for.
\end{itemize}
\end{frame}

% ====================== Slide 6: BACK-OF-ENVELOPE (PIGOVIAN TAX) ======================
\begin{frame}[t,shrink=11]
\frametitle{Back-of-the-Envelope: Appraising Quieter Technology}
\textbf{Illustrative assumptions:} A quieter-technology upgrade costs HK\$3M and reduces expected annual noise charges, monitoring cost, and disruption by HK\$1.2M.
\begin{columns}[T]
\begin{column}{0.48\textwidth}\begin{block}{Without the Upgrade}
\begin{itemize}
\item Expected annual noise-related cost: HK\$1.5M
\item Five-year undiscounted cost: HK\$7.5M
\end{itemize}\end{block}\end{column}
\begin{column}{0.48\textwidth}\begin{block}{With the Upgrade}
\begin{itemize}
\item Upfront cost: HK\$3M
\item Residual annual cost: HK\$0.3M
\item Five-year undiscounted total: HK\$4.5M
\item \textbf{Illustrative five-year saving: HK\$3M}
\end{itemize}\end{block}\end{column}
\end{columns}
\textit{Takeaway: The investment may be worthwhile under these assumptions. A complete appraisal should discount cash flows and test performance, maintenance, and demand effects.}
\vfill{\tiny Illustrative figures; not Hong Kong noise limits, charges, or technology quotations.}
\end{frame}

% ====================== Slide 7: PRIVACY \& DATA ======================
\begin{frame}[t]
\frametitle{Privacy Governance for Onboard Cameras}
Cameras may support navigation, inspection, security, or evidence collection, but privacy risk depends on what is captured, retained, linked, and used.
\begin{itemize}
\item Define a lawful and specific purpose before collecting personal data.
\item Test necessity and proportionality and consider less intrusive alternatives.
\item Minimize field of view, resolution, recording duration, identifiers, audio, and retention where they are not needed.
\item Use access controls, encryption, logs, retention limits, deletion procedures, notices, and incident response.
\item Onboard masking may be one control, but it is not universally required or sufficient by itself.
\end{itemize}
\textit{Hong Kong privacy guidance for drone cameras emphasizes avoiding excessive collection, notification, security, and responsible data handling.}
\end{frame}

% ====================== NEW SLIDE 7.5: PRIVACY EXTERNALITY & MASKING COST ======================
\begin{frame}[t]
\frametitle{Applied Tool: Privacy-Control Cost}
\begin{block}{Annual Privacy-Control Cost}
\[
\text{Annual Cost}=\text{Fixed Governance and System Cost}
+q\times\text{Variable Processing Cost}
+\text{Audit and Incident Cost}
\]
\end{block}
\begin{itemize}
\item Fixed cost may include privacy-impact assessment, system design, notices, access controls, training, and vendor assurance.
\item Variable cost may include secure processing, storage, communications, review, and deletion.
\item The cost cannot be inferred from an unsupported percentage of flight operating cost.
\item Compare controls using privacy effectiveness, safety compatibility, reliability, and total cost.
\end{itemize}
\end{frame}


% ====================== Slide 8: BACK-OF-THE-ENVELOPE - PRIVACY MASKING COST ======================
\begin{frame}[t,shrink=11]
\frametitle{Back-of-the-Envelope: Privacy-Control Cost}
\textbf{Illustrative assumptions:} 50,000 operations/year; fixed privacy-system and governance cost HK\$600k/year.
\begin{columns}[T]
\begin{column}{0.48\textwidth}\begin{block}{Lower Variable Cost}
\begin{itemize}
\item Processing cost: HK\$4/operation
\item Variable cost: HK\$200k
\item Fixed cost: HK\$600k
\item \textbf{Annual cost: HK\$800k}
\end{itemize}\end{block}\end{column}
\begin{column}{0.48\textwidth}\begin{block}{Higher-Assurance Design}
\begin{itemize}
\item Processing cost: HK\$8/operation
\item Variable cost: HK\$400k
\item Fixed cost: HK\$1M
\item \textbf{Annual cost: HK\$1.4M}
\end{itemize}\end{block}\end{column}
\end{columns}
\textit{Takeaway: The more expensive design is justified only if its additional controls materially improve privacy, security, or operational assurance.}
\vfill{\tiny Illustrative assumptions; not a compliance quotation or claim that either design is legally sufficient.}
\end{frame}

% ====================== NEW SLIDE 8.5: CARBON EXTERNALITY - PIGOVIAN CARBON TAX ======================
\begin{frame}[t,shrink=11]
\frametitle{Back-of-the-Envelope: Carbon-Cost Scenario}
\textbf{Illustrative assumptions:} 50,000 operations/year and a carbon value of HK\$1 per kg $\mathrm{CO_2e}$.
\begin{columns}[T]
\begin{column}{0.48\textwidth}\begin{block}{40 kg $\mathrm{CO_2e}$/Operation}
\begin{itemize}
\item Emissions: $50{,}000\times40=2{,}000{,}000$ kg
\item \textbf{Annual carbon cost: HK\$2M}
\end{itemize}\end{block}\end{column}
\begin{column}{0.48\textwidth}\begin{block}{10 kg $\mathrm{CO_2e}$/Operation}
\begin{itemize}
\item Emissions: $50{,}000\times10=500{,}000$ kg
\item \textbf{Annual carbon cost: HK\$0.5M}
\end{itemize}\end{block}\end{column}
\end{columns}
\textit{Takeaway: A carbon value can support investment appraisal, but this is not a claim that CAD imposes such a tax. Emissions factors and boundaries must be sourced.}
\vfill{\tiny Illustrative carbon price and emissions intensities.}
\end{frame}


% ====================== Slide 8: SESSION 2 TITLE ======================
\begin{frame}[c]
  \begin{center}
    \huge\color{cuhkpurple}\textbf{Session 2}\\
    \vspace{0.5cm}
    \Large Carbon Accounting, Life-Cycle Analysis \& Public Acceptance
  \end{center}
\end{frame}

% ====================== Slide 9: CONCEPT PRIMER 2 - LIFE-CYCLE CARBON ACCOUNTING ======================
\begin{frame}[t]
\frametitle{Concept Primer: GHG Inventory and Life-Cycle Assessment}
Corporate GHG scopes and product life-cycle assessment answer related but different questions.
\begin{itemize}
\item \textbf{Scope 1}: Direct emissions from sources owned or controlled by the reporting organization.
\item \textbf{Scope 2}: Indirect emissions associated with purchased electricity, steam, heat, or cooling.
\item \textbf{Scope 3}: Other indirect value-chain emissions, potentially including aircraft, batteries, infrastructure, maintenance, and end-of-life activities.
\item \textbf{Life-cycle assessment}: Evaluates environmental impacts of a defined product or service across a specified life cycle and functional unit.
\end{itemize}
\textit{An electric aircraft may have zero in-flight exhaust emissions while still having Scope 2 and Scope 3 emissions. Corporate scopes should not simply be added without checking boundaries and double counting.}
\end{frame}

% ====================== Slide 10: APPLIED TOOL - GRID CARBON INTENSITY (Scope 2) ======================
\begin{frame}[t]
\frametitle{Applied Tool: Electricity-Related Emissions}
\begin{block}{Operational Electricity Emissions}
\[
\mathrm{CO_2e\ per\ operation}
=\text{kWh consumed}\times\text{electricity emissions factor}
\]
\end{block}
\begin{itemize}
\item Match units carefully, such as kg $\mathrm{CO_2e}$/kWh.
\item Include charging losses and relevant facility energy if they fall within the boundary.
\item Distinguish location-based grid factors from market-based contractual claims where applicable.
\item Compare modes using a consistent functional unit, such as passenger-km, tonne-km, or completed mission.
\end{itemize}
\end{frame}

% ====================== Slide 11: BACK-OF-THE-ENVELOPE - DIRTY GRID ======================
\begin{frame}[t,shrink=10]
\frametitle{Back-of-the-Envelope: Electricity Emissions}
\textbf{Assumption:} 50 kWh of electricity per operation, including charging losses only if stated.
\begin{columns}[T]
\begin{column}{0.48\textwidth}\begin{block}{Factor: 0.8 kg/kWh}
\begin{itemize}
\item $50\times0.8$
\item \textbf{40 kg $\mathrm{CO_2e}$/operation}
\end{itemize}\end{block}\end{column}
\begin{column}{0.48\textwidth}\begin{block}{Factor: 0.2 kg/kWh}
\begin{itemize}
\item $50\times0.2$
\item \textbf{10 kg $\mathrm{CO_2e}$/operation}
\end{itemize}\end{block}\end{column}
\end{columns}
\textit{Takeaway: Electricity source matters, but these figures alone do not establish whether the service outperforms another transport mode. Occupancy, distance, infrastructure, and life-cycle emissions also matter.}
\vfill{\tiny Illustrative factors; use current, sourced factors for an actual inventory.}
\end{frame}

% ====================== NEW SLIDE 12: APPLIED TOOL - FULL LIFE-CYCLE ANALYSIS ======================
\begin{frame}[t]
\frametitle{Applied Tool: Life-Cycle Emissions per Functional Unit}
\begin{block}{Simplified Allocation}
\[
\text{Life-Cycle Emissions per Unit}
=
\frac{E_{vehicle}+E_{battery}+E_{infrastructure}+E_{operations}+E_{end\text{-}of\text{-}life}}
{\text{Lifetime Functional Units}}
\]
\end{block}
\begin{itemize}
\item Functional units may be passenger-km, tonne-km, or completed service missions.
\item State lifetime, utilization, replacement cycles, occupancy or payload, system boundary, and allocation method.
\item Scope 3 is not a universal fixed percentage of total emissions.
\item Compare alternatives using the same boundary and functional unit.
\end{itemize}
\end{frame}

% ====================== NEW SLIDE 13: BACK-OF-THE-ENVELOPE - FULL LCA ======================
\begin{frame}[t,shrink=11]
\frametitle{Back-of-the-Envelope: Allocating Embodied Emissions}
\textbf{Illustrative assumptions:} Embodied aircraft, battery, and infrastructure emissions allocated at 12 kg $\mathrm{CO_2e}$ per operation.
\begin{columns}[T]
\begin{column}{0.48\textwidth}\begin{block}{Higher-Carbon Electricity}
\begin{itemize}
\item Electricity emissions: 40 kg
\item Allocated embodied emissions: 12 kg
\item \textbf{Total: 52 kg $\mathrm{CO_2e}$/operation}
\end{itemize}\end{block}\end{column}
\begin{column}{0.48\textwidth}\begin{block}{Lower-Carbon Electricity}
\begin{itemize}
\item Electricity emissions: 10 kg
\item Allocated embodied emissions: 12 kg
\item \textbf{Total: 22 kg $\mathrm{CO_2e}$/operation}
\end{itemize}\end{block}\end{column}
\end{columns}
\textit{Takeaway: Decarbonizing electricity reduces operational emissions, while embodied emissions remain sensitive to lifetime use and replacement assumptions.}
\vfill{\tiny Illustrative allocation; not a comparative claim about jets or sustainable aviation fuel.}
\end{frame}

% ====================== Slide 14: CONCEPT PRIMER 3 - AiRMOUR PUBLIC ACCEPTANCE ======================
\begin{frame}[t]
\frametitle{Concept Primer: Public Acceptance and Use Cases}
AiRMOUR examined sustainable UAM through the lens of emergency medical services while also considering other applications.
\begin{itemize}
\item Acceptance may vary with perceived public benefit, urgency, safety, noise, privacy, fairness, trust, location, and available alternatives.
\item Urgent medical applications may receive stronger support than some non-urgent applications, but acceptance is not automatic or perfectly divided into two groups.
\item Infrastructure maintenance, surveying, and environmental monitoring may also be viewed favourably in some contexts.
\item Local engagement and evidence are necessary before transferring European findings to Hong Kong.
\end{itemize}
\textbf{Managerial lesson:} Segment by use case and affected stakeholder, then test assumptions rather than assigning universal acceptance scores.
\end{frame}

% ====================== Slide 13: SESSION 3 TITLE ======================
\begin{frame}[c]
  \begin{center}
    \huge\color{cuhkpurple}\textbf{Session 3}\\
    \vspace{0.5cm}
    \Large Equity, Social License \& In-Class Exercise
  \end{center}
\end{frame}

% ====================== Slide 14: CONCEPT PRIMER 4 ======================
\begin{frame}[t]
\frametitle{Concept Primer: Equity and Social Licence}
Social licence refers to continuing acceptance and trust among affected communities and stakeholders. It complements, but does not replace, legal approval.
\begin{itemize}
\item Ask who receives time savings, service improvements, employment, and business opportunities.
\item Ask who bears noise, privacy, safety, land-use, and fiscal costs.
\item Consider affordability, accessibility, geographic coverage, consultation, complaints, and remedies.
\item Acceptance may change over time and cannot be purchased by one symbolic public-service commitment.
\end{itemize}
\textbf{Managerial lesson:} Build legitimacy through performance, transparency, participation, mitigation, and fair distribution of benefits and costs.
\end{frame}

% ====================== Slide 15: STANDALONE FORMULA ======================
\begin{frame}[t]
\frametitle{Applied Tool: Funding a Public-Service Commitment}
\begin{block}{Net Contribution}
\[
\text{Net Contribution}
=Q_C m_C+\text{Public-Service Payment}
-Q_P c_P-\text{Fixed Cost}
\]
\end{block}
\begin{itemize}
\item $Q_Cm_C$ is contribution from commercial services.
\item $Q_Pc_P$ is the incremental cost of the public-service commitment.
\item Funding may come from commercial cross-subsidy, public procurement, availability payment, philanthropy, or another transparent source.
\item The arrangement should specify eligibility, service level, opportunity cost, accountability, and long-term affordability.
\end{itemize}
\end{frame}

% ====================== Slide 16: BACK-OF-ENVELOPE (EQUITY) ======================
\begin{frame}[t,shrink=11]
\frametitle{Back-of-the-Envelope: Cost of a Public-Service Commitment}
\textbf{Assumptions:} 100 daily slots; commercial contribution HK\$400/flight; public-service incremental cost HK\$100/flight.
\begin{columns}[T]
\begin{column}{0.48\textwidth}\begin{block}{All Commercial Slots}
\begin{itemize}
\item 100 commercial flights
\item \textbf{Contribution before fixed cost: HK\$40,000/day}
\item No specified public-service capacity
\end{itemize}\end{block}\end{column}
\begin{column}{0.48\textwidth}\begin{block}{Twenty Priority Slots}
\begin{itemize}
\item Commercial contribution: $80\times400=\text{HK\$32,000}$
\item Public-service cost: $20\times100=\text{HK\$2,000}$
\item \textbf{Net before fixed cost: HK\$30,000/day}
\item Opportunity cost relative to all-commercial case: HK\$10,000/day
\end{itemize}\end{block}\end{column}
\end{columns}
\textit{Takeaway: The commitment has a measurable cost. Whether it is justified and who should fund it require a public-value and distributional assessment.}
\vfill{\tiny Illustrative assumptions; no claim that the commitment secures approval or acceptance.}
\end{frame}

% ====================== Slide 17: APPLIED TOOL - SUSTAINABILITY KPIs ======================
\begin{frame}[t]
\frametitle{Applied Tool: Sustainability KPI Design}
\begin{center}\small
\begin{tabular}{l|l|l}
\textbf{Area} & \textbf{Example KPI} & \textbf{Boundary or denominator}\\
\hline
Noise & Population-weighted exposure & Route, period, and noise metric\\
Climate & kg $\mathrm{CO_2e}$/passenger-km & Life-cycle boundary and occupancy\\
Privacy & Personal-data incidents & Operations and severity\\
Equity & Service availability by area & Population, need, and affordability\\
Governance & Complaints resolved on time & Defined process and deadline\\
\end{tabular}
\end{center}
\vspace{0.25cm}
\begin{itemize}
\item Define baseline, target, measurement method, frequency, owner, assurance, and corrective action.
\item No universal CAD or investor requirement mandates the three ratios previously shown.
\item A KPI tracks performance; it does not by itself prove profitability or social acceptance.
\end{itemize}
\end{frame}

% ====================== Slide 18: EXERCISE INSTRUCTIONS ======================
\begin{frame}[t]
\frametitle{In-Class Exercise: Sustainability Agreement}
\textbf{Scenario:} A high-volume corridor is proposed near a residential area.
\begin{itemize}
\item \textbf{Roles:} Operator, community representatives, regulator, public-service customer, and data-protection adviser
\item \textbf{Tasks:}
\begin{enumerate}
\item identify noise, privacy, climate, access, and distributional impacts;
\item propose one prevention or mitigation measure for each material impact;
\item estimate cost, effectiveness, beneficiary, payer, and residual impact;
\item define five measurable KPIs and review triggers; and
\item agree on complaints, incident, pause, and corrective-action procedures.
\end{enumerate}
\item \textbf{Timing:} 15 minutes plus 5-minute discussion.
\end{itemize}
\end{frame}

% ====================== Slide 19: EXERCISE TEMPLATE ======================
\begin{frame}[t]
\frametitle{Exercise Template: Mitigation Scorecard}
\begin{center}\small
\begin{tabular}{l|l|l|l}
\textbf{Impact} & \textbf{Measure} & \textbf{KPI} & \textbf{Payer}\\
\hline
Noise & Route, schedule, or technology & Exposure metric & To negotiate\\
Privacy & Data minimization and security & Incident and retention metric & To negotiate\\
Climate & Energy and utilization strategy & kg $\mathrm{CO_2e}$/functional unit & To negotiate\\
Equity & Access or public-service design & Availability and affordability & To negotiate\\
\end{tabular}
\end{center}
\vspace{0.3cm}
\textit{Do not assume a 60 dB limit or a fixed number of free public-service flights. Define thresholds from evidence, regulation, operating context, and stakeholder agreement.}
\end{frame}

% ====================== Slide 20: KEY TAKEAWAYS ======================
\begin{frame}[t]
\frametitle{Key Takeaways}
\begin{itemize}
\item Externalities should be measured and addressed through a mix of technology, standards, operating rules, pricing, compensation, and engagement.
\item Noise cannot be priced reliably by multiplying flights by decibels above an assumed threshold.
\item Privacy governance requires lawful purpose, necessity, proportionality, minimization, security, transparency, retention control, and accountability.
\item Corporate GHG scopes and life-cycle assessment require clear and consistent boundaries; electric operation does not imply zero life-cycle emissions.
\item Public acceptance varies across use cases and stakeholders and cannot be reduced to a universal medical-versus-commercial split.
\item Equity commitments and KPIs need measurable outcomes, transparent funding, review, and corrective action.
\end{itemize}
\end{frame}

% ====================== Slide 21: ASSIGNMENT ======================
\begin{frame}[t]
\frametitle{Week 10 Assignment Briefing}
\textbf{Sustainability and Social-Impact Plan (maximum 1,200 words plus KPI appendix)}
\begin{itemize}
\item \textbf{Due:} Beginning of Week 11 (group assignment)
\item \textbf{Requirements:}
\begin{enumerate}
\item define the principal noise, privacy, climate, and distributional impacts;
\item propose a mitigation hierarchy: avoid, reduce, manage, and compensate residual impacts;
\item calculate operational electricity emissions and one life-cycle indicator with sourced factors;
\item prepare a privacy-governance plan covering necessity, minimization, security, retention, notice, and incidents;
\item evaluate one equity or public-service commitment, including cost, beneficiary, and payer; and
\item define at least five KPIs with boundaries, baselines, targets, measurement methods, owners, and review triggers.
\end{enumerate}
\item \textbf{Grading focus:} Economic reasoning (25\%), environmental analysis (25\%), privacy and governance (20\%), equity and stakeholder analysis (20\%), professionalism (10\%).
\end{itemize}
\end{frame}

% ====================== Slide 22: PREVIEW WEEK 11 ======================
\begin{frame}[t]
  \frametitle{Preview of Week 11 \& Q\&A}
  
  \begin{center}
      \Large \textbf{Synthesis \& Capstone Group Presentations}
  \end{center}
  
  \vspace{0.5cm}
  \textbf{Preview question for next week:} \\
  \textit{How will you combine utility formulas, unit economics, infrastructure constraints, and sustainability into one master presentation that secures \$100M in venture funding?}
  
  \vspace{0.5cm}
  \begin{center}
      \textbf{Questions?} \\
      \small (Final 15 minutes reserved for extended Q\&A)
  \end{center}
\end{frame}

\end{document}